\begin{theorem} \label{teo:fubini_r3} \textbf{Teorema de Fubini} en $\mathbb{R}^3.$
    Sea $f:W\to\mathbb{R}$ un campo escalar continuo sobre $W\subseteq\mathbb{R}^3$ un conjunto elemental.
    
    \begin{enumerate} 
        \item Si $W$ es un conjunto de \textbf{Tipo 1}  (\ref{def_conj_elem_R3_t1}), entonces la integral triple de $f$ sobre $W$ se puede calcular como:
        \[
            \iiint_W f(x,y,z) \, dV = \iint_{D_1} \left( \int_{\eta_{11}(x,y)}^{\eta_{21}(x,y)} f(x,y,z) \, dz \right) dA
        \]
       
        \item Si $W$ es un conjunto de \textbf{Tipo 2} (\ref{def_conj_elem_R3_t2}), entonces la integral triple de $f$ sobre $W$ se puede calcular como:
        \[
            \iiint_W f(x,y,z) \, dV = \iint_{D_2} \left( \int_{\eta_{12}(x,z)}^{\eta_{22}(x,z)} f(x,y,z) \, dy \right) dA
        \]
       
 \item Si $W$ es un conjunto de \textbf{Tipo 3} (\ref{def_conj_elem_R3_t3}), entonces la integral triple de $f$ sobre $W$ se puede calcular como:
        \[
            \iiint_W f(x,y,z) \, dV = \iint_{D_3} \left( \int_{\eta_{13}(y,z)}^{\eta_{23}(y,z)} f(x,y,z) \, dx \right) dA
        \]
    \end{enumerate}
\end{theorem}


